\section*{Bab \ref{ch:g8:fronts-and-weather-maps} --- Massa Udara, Front, dan Peta Cuaca}

\begin{solution}{exo:g8:fronts-and-weather-maps:1}
(a) Maritim tropis: hangat dan lembap. (b) Kontinental polar: dingin dan
kering. (c) Kontinental tropis: panas dan kering. (d) Maritim polar:
dingin dan lembap.
\end{solution}

\begin{solution}{exo:g8:fronts-and-weather-maps:2}
Garis biru bersegitiga: \omterm{def:g8:fronts-and-weather-maps:front}{front dingin}. Garis merah berhiaskan setengah
lingkaran: \omterm{def:g8:fronts-and-weather-maps:front}{front hangat}. Garis yang menghubungkan tempat-tempat yang
tekanannya sama: \omterm{def:g8:fronts-and-weather-maps:isobar}{isobar}. T: pusat \omterm{def:g8:fronts-and-weather-maps:depression}{antisiklon} (tekanan tinggi). R: pusat
\omterm{def:g8:fronts-and-weather-maps:depression}{depresi} (tekanan rendah).
\end{solution}

\begin{solution}{exo:g8:fronts-and-weather-maps:3}
(a) Benar. (b) Salah: \omterm{def:g4:clouds-rain-snow:cloud}{awannya} datang sebelum \omterm{def:g8:fronts-and-weather-maps:front}{front}, jauh di depannya,
tinggi di langit. (c) Salah: \omterm{def:g8:fronts-and-weather-maps:isobar}{isobar} yang berdekatan berarti angin yang
kencang. (d) Benar.
\end{solution}

\begin{solution}{exo:g8:fronts-and-weather-maps:4}
(a) \qty{980}{hPa}. (b) \qty{1034}{hPa}. (c) \qty{4}{hPa}. (d) $1034 -
980 = \qty{54}{hPa}$.
\end{solution}

\begin{solution}{exo:g8:fronts-and-weather-maps:5}
Guratan tipis \omterm{def:g4:clouds-rain-snow:cloud}{awan} \omterm{def:g4:clouds-rain-snow:cloud-families}{sirus}; selubung kelabu altostratus; \omterm{def:g4:clouds-rain-snow:cloud}{awan} yang makin
rendah dan tebal; \omterm{def:g4:clouds-rain-snow:precipitation}{hujan} yang tetap; udara yang lebih hangat dan lebih
lembap.
\end{solution}

\begin{solution}{exo:g8:fronts-and-weather-maps:6}
Pada \omterm{def:g8:fronts-and-weather-maps:front}{front dingin}, udara dingin yang rapat menyusup ke bawah udara hangat
seperti \emph{baji}: udara hangat terangkat pada lereng yang \emph{curam},
naik dengan cepat, dan membentuk \omterm{def:g4:clouds-rain-snow:cloud}{awan} yang tinggi. Pada \omterm{def:g8:fronts-and-weather-maps:front}{front hangat},
udara hangat naik perlahan-lahan di atas udara dingin pada \emph{lereng
landai}: udara itu naik sedikit demi sedikit sepanjang jarak yang jauh,
dan \omterm{def:g4:clouds-rain-snow:cloud}{awannya} menyebar berlapis-lapis.
\end{solution}

\begin{solution}{exo:g8:fronts-and-weather-maps:7}
(a) Di dekat R: di sana \omterm{def:g8:fronts-and-weather-maps:isobar}{isobar-isobar} lebih berdekatan, dan panahnya
lebih panjang. (b) Dari timur (panah-panahnya menunjuk ke barat, sedikit
ke arah barat daya). (c) Dari timur (panah-panahnya menunjuk ke barat).
(d) Berlawanan arah jarum jam.
\end{solution}

\begin{solution}{exo:g8:fronts-and-weather-maps:8}
(a) $1016 - 1003 = \qty{13}{hPa}$. (b) \omterm{def:g8:fronts-and-weather-maps:depression}{Depresi}: tekanannya terus turun.
(c) \omterm{def:g4:clouds-rain-snow:cloud}{Awan}, angin, dan \omterm{def:g4:clouds-rain-snow:precipitation}{hujan}, mungkin disertai sebuah \omterm{def:g8:fronts-and-weather-maps:front}{front}.
\end{solution}

\begin{solution}{exo:g8:fronts-and-weather-maps:9}
(a) cT, kontinental tropis. (b) Panas dan kering. (c) Di atas \omterm{def:g2:water-on-earth:ocean}{laut}, udara
itu menyerap uap air: udara itu menjadi lebih lembap, dan udara di dekat
permukaan air sedikit didinginkan oleh \omterm{def:g2:water-on-earth:ocean}{laut}.
\end{solution}

\begin{solution}{exo:g8:fronts-and-weather-maps:10}
(a) $4 - 1 = 3$ hari. (b) $3 \div 3 = 1$ hari untuk setiap sepuluh tahun.
(c) Diperlukan $5 - 1 = 4$ hari lagi: sekitar 40 tahun.
\end{solution}

\begin{solution}{exo:g8:fronts-and-weather-maps:11}
(a) Sekitar \qty{1350}{km} (empat setengah kali batang \omterm{def:g2:drawing-to-map:scale}{skala}). (b)
\qty{4}{hPa} sepanjang \qty{80}{km}:
$4 \div 80 \times 100 = \qty{5}{hPa}$ per \qty{100}{km}.
(c) $4 \div 140 \times 100 \approx \qty{2.9}{hPa}$ per \qty{100}{km}. (d) Di dekat R, tempat tekanan berubah
lebih cepat; ya, panah-panahnya lebih panjang di dekat R.
\end{solution}

\begin{solution}{exo:g8:fronts-and-weather-maps:12}
(a) \omterm{def:g8:fronts-and-weather-maps:front}{Front dingin}: $300 \div 50 = 6$ jam. \omterm{def:g8:fronts-and-weather-maps:front}{Front hangat} bergerak dengan
kecepatan \qty{25}{km} per jam: $200 \div 25 = 8$ jam. \omterm{def:g8:fronts-and-weather-maps:front}{Front dingin} tiba
2 jam lebih awal. (b) Kota pertama: \omterm{def:g4:clouds-rain-snow:cloud}{awan} yang menjulang, \omterm{def:g4:clouds-rain-snow:precipitation}{hujan} deras,
embusan angin, lalu udara yang lebih sejuk dan langit yang makin cerah.
Kota kedua: \omterm{def:g4:clouds-rain-snow:cloud}{awan} berlapis yang makin rendah dan \omterm{def:g4:clouds-rain-snow:precipitation}{hujan} yang panjang dan
tetap, lalu udara yang lebih hangat dan lebih lembap. (c) Kota kedua,
yang didatangi udara hangat.
\end{solution}

\begin{solution}{pb:g8:fronts-and-weather-maps:1}
\textbf{1.} Di dalam udara hangat, di antara kedua \omterm{def:g8:fronts-and-weather-maps:front}{front}. \omterm{def:g8:fronts-and-weather-maps:front}{Front hangat}
sudah lewat.

\textbf{2.} Dari barat daya.

\textbf{3.} Ke timur: segitiga-segitiga \omterm{def:g8:fronts-and-weather-maps:front}{front dingin} terletak di sisi ke
arah mana \omterm{def:g8:fronts-and-weather-maps:front}{front} itu bergerak, dan segitiga-segitiga itu menghadap ke
timur, ke arah kota.

\textbf{4.} \omterm{def:g4:clouds-rain-snow:cloud}{Awan} yang menjulang, \omterm{def:g4:clouds-rain-snow:precipitation}{hujan} deras, mungkin guntur dan \omterm{def:g4:clouds-rain-snow:precipitation}{hujan es},
embusan angin; lalu udara yang lebih sejuk, angin dari barat atau barat
laut, dan langit yang makin cerah.

\textbf{5.} Di \qty{180}{km} pada pukul 06.00 dan di \qty{300}{km} pada pukul 09.00.

\textbf{6.} $300 - 180 = \qty{120}{km}$.

\textbf{7.} $120 \div 3 = \qty{40}{km}$ per jam.

\textbf{8.} Kota itu berada di \qty{480}{km}: $480 - 300 = \qty{180}{km}$.

\textbf{9.} $180 \div 40 = 4.5$ jam, yaitu 4 jam 30 menit.

\textbf{10.} 09.00 + 4~jam~30~menit = pukul \textbf{13.30}.

\textbf{11.} Tidak: \omterm{def:g4:clouds-rain-snow:precipitation}{hujan} deras seharusnya berhenti sekitar pukul 14.30,
sebelum sepak mula.

\textbf{12.} $180 \div 30 = 6$ jam: pada pukul 15.00, tepat pada saat
sepak mula. Prakiraan: ``\omterm{def:g8:fronts-and-weather-maps:front}{Front dingin} diperkirakan mencapai kota sekitar
pukul \textbf{13.30}, disertai \omterm{def:g4:clouds-rain-snow:precipitation}{hujan} deras dan embusan angin selama satu
jam, lalu langit akan cerah: pertandingan kemungkinan besar dapat
dimainkan. Jika \omterm{def:g8:fronts-and-weather-maps:front}{front} melambat, \omterm{def:g8:fronts-and-weather-maps:front}{front} itu bisa tiba pada pukul 15.00 dan
menunda pertandingan.''
\end{solution}
